\begin{wideslide}{Buffon's needle problem}
\null\vfill

  \twocolumn
  {
    {\it Suppose we have a floor made of parallel strips of wood, each the same width, and we drop a needle onto the floor. What is the probability that the needle will lie across a line between two strips?}
    
    \sep
    
    {\it\hfill Georges-Louis Leclerc,\\\hfill Comte de Buffon\\\hfill 18th century}
  }
  {
    \sep\sep
    \centering\input{figures/buffon.tex}
  }
  
  \vspace{-10pt}
  \myBoxFullWidth{Monte Carlo without computers}
  
  \twocolumn
  {
    If needle length ($l$) $<$ lines width ($t$):
    
    $$P = \frac{2l}{t\pi}$$
    
    which can be used to estimate $\pi$:
    
    $$\pi = \frac{2l}{tP}$$
  }
  {
    MC experiment was performed by Mario Lazzarini in 1901 by throwing 3408 needles:
    
    $$\pi = \frac{2l \cdot 3408}{t \cdot \#red} = \frac{355}{113} = 3.14159292$$
  }
  
\vfill\null
\end{wideslide}

\begin{wideslide}[toc = From Solitaire to MC]{From Solitaire to Monte Carlo method}
\null\vfill

    \twocolumn
    {
      \begin{itemize}
	\item Stanis{\l}aw Ulam was a Polish mathematician
	\item He invented the Monte Carlo method while playing solitaire
	\item The method was used in Los Alamos, performed by ENIAC computer
      \end{itemize}  
      
      \includegraphics[width=\columnwidth]{figures/eniac1946.eps}
    }
    {
      \includegraphics[width=\columnwidth]{figures/solitaire.eps}
      \begin{itemize}
	\item What is a probability of success in solitaire?
	\begin{itemize}
	  \item Too complex for an analytical calculations
	  \item Lets try $N = 100$ times and count wins
	  \item With $N \rightarrow \infty$ we are getting closer to correct result
	\end{itemize}
      \end{itemize}
    }
 
\vfill\null
\end{wideslide}

\begin{wideslide}[toc = Simple Monte Carlo, method=direct]{Simple Monte Carlo - coin flipping}
\null\vfill

  \begin{itemize}
    \item What is the probability of getting head / tail when throwing a coin?
  \end{itemize}
  
  \begin{minipage}{0.7\textwidth}
  \begin{itemize}
    \item One can use MC method to calculate this:
    \begin{itemize}
      \item get $N$ times a random number $x \in [0,1)$
      \item check how many ($n$) times $x < 0.5$
      \item $P = \frac{n}{N}$
    \end{itemize}
    \item One-line example in Python:
  \end{itemize}
  \end{minipage}\begin{minipage}{0.3\textwidth}
		  \includegraphics[width=0.5\columnwidth]{figures/coin.eps}
                \end{minipage}

  {\small\color{pdcolor3}
  \begin{verbatim}
MC = lambda N: len(filter(lambda x: random.random() < 0.5, range(N))) / float(N)     
  \end{verbatim}
  }
  \begin{minipage}{0.55\textwidth}
  \begin{itemize}
   \item If you do not like one-liners:
  \end{itemize}
  \vspace{-20pt}  
  {\small\color{pdcolor3}
  \begin{verbatim}
def MonteCarlo(N):
  n = 0
  for _ in range(N):
    if (random.random() < 0.5): n += 1
  return float(n) / N   
  \end{verbatim}
  }
  \end{minipage}\begin{minipage}{0.45\textwidth}
		  \vspace{-10pt}
		  \begin{itemize}
		   \item Some results for MC(100):
		   \item[]
		   0.51, 0.48, 0.61, 0.46, 0.57, 0.46
		   \item Some results for MC(100000):
		   \item[]
		   0.50139, 0.49798, 0.49933, 0.5005
		  \end{itemize}
                \end{minipage}

\vfill\null
\end{wideslide}